\subsection{\texorpdfstring{\(A^{+}(\mathrm{X})\) 到 \(\mathrm{L}(\mathrm{X})\) 上的投影}{A+ (X) 到 L(X) 上的投影}}

设 \(\pi\) 是 \(A^{+}(X)\) 到 \(L(X)\) 中的由下式定义的线性映射

\[
\pi (x _ {1} \dots x _ {n}) = (\mathrm{ad} (x _ {1}) \circ \dots \circ \mathrm{ad} (x _ {n - 1})) (x _ {n})\tag{3}
\]

其中 X 中的 \(n \geqslant 1, x_1, \ldots, x_n\) 。

命题 1. (a) 限制 \(\pi_0\) （即 \(\pi\) 在 \(\mathrm{L}(\mathbf{X})\) 上的限制）是 \(\mathrm{L}(\mathbf{X})\) 的一个导子。

\begin{enumerate}
\def\labelenumi{(\alph{enumi})}
\setcounter{enumi}{1}
\item
  对每个整数 \(n \geqslant 1\) 以及一切 \(u\) （属于 \(\mathbf{L}^n(\mathbf{X})\) ）， \(\pi(\mathbf{u}) = n.u.\)
\item
  设 E 是模 L(X) 的自同态代数， \(\theta\) 是 A(X) 到 E 中的一个同态，使得 \(\theta(x) = \operatorname{ad} x\) （对一切 \(x \in X\) ）。 \(\theta\) 在 L(X) 上的限制是 L(X) 到 E 中的一个李代数同态，它在 X 上与 L(X) 的伴随表示一致，从而
\end{enumerate}

\[
\theta (u). v = [ u, v ] \quad \text { for   } u, v \text {   in   } L (X).\tag{4}
\]

设 \(a\) 属于 \(\mathrm{A}(\mathrm{X})\) ， \(b\) 属于 \(\mathrm{A}^{+}(\mathrm{X})\) ；则

\[
\pi (a. b) = \theta (a). \pi (b).\tag{5}
\]

只需考虑 \(a = x_{1} \ldots x_{p}\) 、 \(b = x_{p+1} \ldots x_{p+q}\) （其中 \(p \geqslant 0\) ）、 \(q \geqslant 1\) 以及 \(x_{1}, \ldots, x_{p+q}\) （属于 \(\mathbf{X}\) ）的情形；但这时 (5) 由 (3) 立即得出，因为 \(\theta(x) = \operatorname{ad} x\) （对 \(x \in \mathbf{X}\) ）。

设 \(u\) 与 \(v\) 属于 \(\mathbf{L}(\mathbf{X})\) ；由 (4) 与 (5)，

\[
\begin{array}{r l} \pi_ {0} ([ u, v ]) & = \pi (u v - v u) = \theta (u). \pi (v) - \theta (v). \pi (u) \\ & = [ u, \pi (v) ] - [ v, \pi (u) ] = [ u, \pi_ {0} (v) ] + [ \pi_ {0} (u), v ], \end{array}
\]

因此 \(\pi_0\) 是 L(X) 的一个导子。

\begin{enumerate}
\def\labelenumi{(\alph{enumi})}
\setcounter{enumi}{1}
\tightlist
\item
  设 \(\pi_1\) 是模 \(\mathrm{L}(\mathbf{X})\) 的这样一个自同态：它在 \(\mathrm{L}^n (\mathbf{X})\) 上与用整数 \(n\geqslant 1\) 作乘法一致。公式 \([\mathrm{L}^n (\mathbf{X}),\mathrm{L}^m (\mathbf{X})]\subset \mathrm{L}^{n + m}(\mathrm{X})\) 表明 \(\pi_{1}\) 是一个导子 (《代数》， 第 III 章， § 10， no. 3， 例 6)。导子 \(\pi_1 - \pi_0\) （\(\mathrm{L}(\mathbf{X})\) 的）在 \(\mathbf{X}\) 上为零，并且由于 \(\mathbf{X}\) 生成 \(\mathrm{L}(\mathbf{X}),\pi_0 = \pi_1\) ，由此得 (b)。
\end{enumerate}

推论. 假设 K 是一个 Q-代数。设 P 是 \(A^{+}(X)\) 到其自身的线性映射，使得

\[
\mathrm{P} \left(x _ {1} \dots x _ {n}\right) = \frac {1}{n} (\operatorname{ad} x _ {1} \circ \dots \circ \operatorname{ad} x _ {n - 1}) (x _ {n})\tag{6}
\]

其中 \(n \geqslant 1\) 与 \(x_1, \ldots, x_n\) 都属于 \(\mathbf{X}\) 。那么 \(\mathbf{P}\) 是 \(\mathbf{A}^+(\mathbf{X})\) 到 \(\mathbf{L}(\mathbf{X})\) 上的一个投影。

\(\mathbf{P}\) 的像包含在 \(\mathrm{L}(\mathrm{X})\) 中。此外，对一切 \(n \geqslant 1\) 以及一切 \(u\) （属于 \(\mathrm{L}^n (\mathrm{X})\) ），有 \(\mathrm{P}(u) = \frac{1}{n}\pi (u)\) ，从而由命题 1 得 \(\mathrm{P}(u) = u\) 。由于

\[
\mathrm{L} (\mathrm{X}) = \sum_ {n \geqslant 1} \mathrm{L} ^ {n} (\mathrm{X}),
\]

我们看出 P 在 L(X) 上的限制是恒同映射。

注. 假设 K 是一个特征为零的域，并设 Q 是与 U(L(X)) 的典范分次相伴的、A(X) = U(L(X)) 到 L(X) 上的投影，参见 § 1 no. 5。对 \(\alpha \in \mathbf{N}^{(\mathbf{X})}\) ，有 Q(A \(^\alpha\) (X)) ⊂ L \(^\alpha\) (X)。因为只需验证 Q 的像与核都是 A(X) 的带有 N 型 \(^\alpha\) 分次的分次子模。这对像是显然的，因为它等于 L(X)。另一方面，设 n 是一个整数 \(\geqslant\) 1。A(X) 的由诸 \(y^n\) （其中 \(y \in L\) (X)）生成的向量子空间等于 A(X) 的由诸 \(\sum_{\sigma \in \mathcal{E}_n} y_{\sigma(1)} y_{\sigma(2)} \ldots y_{\sigma(n)}\) （其中 \(y_1, y_2, \ldots, y_n\) 是 L(X) 的齐次元素）生成的向量子空间；于是这个子空间是 A(X) 的一个分次子模。

（注意，若 \(\operatorname{Card}(\mathbf{X}) \geqslant 2\) ，则投影 \(\mathbf{P}\) 与 \(\mathbf{Q}\) 在 \(\mathbf{A}^{+}(\mathbf{X})\) 上并不一致。因为设 \(x, y\) 属于 \(\mathbf{X}\) 且 \(x \neq y\) ，并写

\[
z = x [ x, y ] + [ x, y ] x = x ^ {2} y - y x ^ {2}.
\]

则 \(\mathbf{Q}(z) = 0\) 且 \(\mathrm{P}(z) = \frac{1}{3}[x, [x,y]] \neq 0\) ，参见 § 2 no. 10 的例以及 no. 11 的定理 1。）

